\subsection{Predictions}
\label{sec:design-hypotheses}

The measurements were designed around five predictions. The first four are the statements of the measurements on the default definition, informative withdrawals and thin grades; each is stated with the quantity that decides it and the outcome that would refute it. The fifth is the prediction of the theory. Of the five, only part of H1 (its bootstrap interval), H2, H3 (at its stated level) and H5 are decided by computations that follow their statement; the point values of H1 and all of H4 are numbers that the case had already computed, so that no outcome could have refuted them and they are reported as reproductions. The statement from which the first four are taken reads: ``the same agency's long-run PD of the CCC to C grade is 26.64\% under the CEREP default page and 8.61\% under the transition page, and the IRB risk weight of the cohort mix moves from 58.0\% to 47.5\% with the definition; informative withdrawals and thin grades are measured on rating paths of known truth (with withdrawals ten times as likely in the year before default, removing them understates the PD of B by 34\% and following them as EBA paragraph 76 asks by 15\%, and no matrix test sees it; with 50 obligors per grade the Wald interval contains the PD of AAA in 0.33\% of the samples)''.
\begin{description}
\item[H1, the definition.] ``The same agency's long-run PD of the CCC to C grade is $\spLraDtwoCcc\%$ under the CEREP default page and $\spLraDfourCcc\%$ under the transition page, and the IRB risk weight of the cohort mix moves from $\spRwDtwo\%$ to $\spRwDfour\%$ with the definition.'' The deciding quantities are the gap index $\rho_{\mathrm{CCC}}$ of Equation~\eqref{eq:gap-index}, predicted below $0.5$ with the upper end of its 95\% bootstrap interval below $0.5$, and the fall of the risk weight, predicted above 5 points with the lower end of its interval above zero. The prediction is refuted if either interval reaches its limit. The point values of the two quantities are those of the case's committed computation and are reproduced; the interval is the new computation that decides.
\item[H2, informative withdrawals.] With withdrawal ten times as likely in the year before default ($k=9$), removing the withdrawn ratings understates the PD of B and following them as the guidelines ask does as well, by $\committedremovedkNine\%$ and $\committedfollowedkNine\%$ in the first measurement of 200 repetitions. The deciding quantity is the relative bias of Equation~\eqref{eq:rel-bias} of each treatment in the extension of $\repsNew$ repetitions on new seeds: the prediction is that the upper end of its 95\% interval (the bias plus 1.96 standard errors) lies below $-10\%$ for both treatments. It is refuted if either interval reaches $-10\%$. % voice-ok: withdrawn rating, the domain term of the study
\item[H3, the transition tests.] ``No matrix test sees it.'' The deciding quantity of the first run is the rejection rate at the 5\% level of the test of time homogeneity and of the test of the last matrix against the true one, at $k=9$ over $\repsTests$ repetitions: the prediction is that the upper end of the Wilson interval of each is below 10\%, twice the size of the test. It is refuted if either reaches 10\%. The statement is wider than that test, because the engine has further transition tests; three of them were run afterwards, once the first results were known, on the same draws and at the same three levels of $k$ (the tests of the Markov order, of rating momentum and of the monotonicity of the migration matrix; the matrix weighted bandwidth is a descriptive metric with no rejection rule and is not run). The statement as made is refuted by a rejection rate whose Wilson interval lies above 10\% for any of the five tests at any of the three levels of $k$.
\item[H4, thin grades (a reproduction).] With 50 obligors per grade the Wald interval contains the PD of AAA in $\waldAaaFifty\%$ of the samples, and the Agresti-Coull and Jeffreys intervals cover at least 95\% in every grade at that size. The coverage is computed exactly and the case had already committed it, so the statement is an equality to the stated digits that no outcome could refute: it is reproduced here and not tested.
\item[H5, the closed forms.] H5a: the closed forms of Proposition~\ref{prop:first-year} agree with the means of the engine's simulation of the first cohort of a two-state chain within three Monte Carlo standard errors, for every estimator and every level of $k$. H5b: the numerical evaluation of the model's masses over the five cohorts, which includes the truncation of the window at the horizon of the generator and for which this paper derives no closed form beyond the later-cohort value of Proposition~\ref{prop:later}, agrees with the pooled means within three Monte Carlo standard errors. Either is refuted by a standardised difference above three in absolute value.
\end{description}
The first-order formula of Equation~\eqref{eq:first-order} is stated as a conjecture with its refutation condition in Section~\ref{sec:theory-conjecture} and is judged against the first-cohort measurements in Section~\ref{sec:results-withdrawals}.

The numbers of H1 to H4 are those of the committed computations of the case, which preceded the extension; the extension on new seeds is what tests H2 and H3 at a larger number of repetitions. The thresholds of the table of outcomes were fixed with those numbers in view, a limitation discussed in Section~\ref{sec:limitations}.

\subsection{Data and provenance}
\label{sec:design-data}

The real data are ESMA's CEREP statistics for the EU entities of three agencies, corporate long-term ratings, mapped to seven performing grades (AAA, AA, A, BBB, BB, B, and CCC to C), default and withdrawals, for annual cohorts from 2000 to 2025 for S\&P and Moody's and from 2002 to 2025 for Fitch \citep{ESMA_CEREP}. Table~\ref{tab:data} gives the cohorts, the mean cohort size and the pooled counts of the speculative grade that carries most of the gap. The repository's terms authorise reproduction provided the source is acknowledged, and every table derived from it here carries the attribution that the tables are transformed from ESMA CEREP. The pages differ in more than the definition: the cohort that the default page divides by exceeds the cohort of the transition page in $\spGapYears$ of the $\spCohorts$ years of S\&P, in $\moGapYears$ of the $\moCohorts$ of Moody's and in $\fiGapYears$ of the $\fiCohorts$ of Fitch, and the repository gives no reason. Moody's transition page has no default category, so the transition-page and kept definitions exist for two agencies only. Fitch's transition pages from \fiPopDropped{} have empty default categories, so its transition-page rates in those \fiPopDroppedN{} years are zero where its default page counts defaults; the Fitch results are therefore given on all \fiCohorts{} years and on the \fiPopYears{} years in which the default category is populated. No real data say whether withdrawals were informative; that is why the known-truth experiments exist.

\begin{table}[htbp]
\centering
\setlength{\tabcolsep}{3.5pt}
\footnotesize
\caption{The CEREP data behind the real-data comparison: three agencies' corporate long-term ratings, annual cohorts, the mean number of ratings in a cohort and the pooled defaults over rated ratings of the CCC to C grade under the default page and the transition page.}
\label{tab:data}
\resizebox{\linewidth}{!}{\begin{tabular}{lrlrrr}
\toprule
Agency & Cohorts & Years & Mean cohort size & CCC to C, default page (defaults/ratings) & CCC to C, transition page (defaults/ratings) \\
\midrule
\relax S\&P & \relax 26 & \relax 2000 to 2025 & \relax 7111 & \relax 1489/5428 & \relax 301/3731 \\
\relax Moody's & \relax 26 & \relax 2000 to 2025 & \relax 3758 & \relax 821/6983 & \relax none \\
\relax Fitch & \relax 24 & \relax 2002 to 2025 & \relax 2857 & \relax 331/1553 & \relax 43/1066 \\
\bottomrule
\end{tabular}}
\end{table}


The known-truth data come from the generator of the case. Its latent chain is the expectation-maximisation generator of the annual counts of S\&P, pooled over the cohorts, with seven performing grades and default \citep{dosReis_2018,Bladt_2005}; its one-year PDs run from $\pAaaBp$ basis points for AAA, through $\pAaBp$ basis points for AA (the fitted generator puts AA below AAA, so that the grades order by rating but not by PD at the top), to $\pCcc\%$ for CCC to C, with $\pB\%$ for grade B, and the cohorts have the sizes of S\&P's mean cohort (\genSizes{} obligors from AAA to CCC to C). Obligors are drawn exactly, with annual snapshots over five years, and withdrawn at the rate $w=\wRate\%$ a year with the window of one year of Equation~\eqref{eq:withdrawal-hazard}. The truth status of these data is synthetic known truth, and the results transfer to real data only as far as the generator resembles them (Section~\ref{sec:limitations}). % voice-ok: withdrawn rating, the domain term of the study

\subsection{The methods compared, and how they are run}
\label{sec:design-methods}

Every method compared is run by the product on the same data; no number is taken from a publication. On the real cohorts the four definitions of Equation~\eqref{eq:definitions} are computed from the counts. On the paired draws of the generator the four estimators of Equation~\eqref{eq:estimators} are computed by the pooled cohort estimator of the engine \citep{Anderson_1957}: the removal treatment is the convention that Schuermann and Hanson call the industry standard and that CEREP's transition page implements \citep{Schuermann_2004}; the keeping treatment is the transition page with the withdrawals as an absorbing state; the followed treatment is the reading of paragraphs 73 and 76 of the guidelines \citep{EBA_2017}, which the product implements at the one-year level, reading the end of each rating of the cohort on its latent path; and the no-withdrawal reference follows every obligor. The published estimator that handles a withdrawal as a censoring, the duration estimator of Lando and Skodeberg, is run on the observed continuous-time paths of the same draws: the engine's maximum-likelihood generator of the paths, in which a withdrawal ends the exposure of a rating without a jump, and the default column of its matrix exponential is the one-year PD \citep{Lando_2002}. The engine also provides the Aalen-Johansen estimator of the transition matrix; it is not run, because, like the duration estimator, it treats the censoring as independent of default, and the measurement of the duration estimator already tests that assumption. The interval estimators compared are the Wald, Agresti-Coull and Jeffreys intervals of the engine \citep{Brown_2001,Agresti_1998}, and the screening tests are the engine's tests of time homogeneity and of a matrix against a reference, with the tests of the Markov order, of rating momentum and of the monotonicity of the migration matrix. The case does not run a repair estimator, because the protocol is a measurement.

\subsection{Metrics, statistics and protocol}
\label{sec:design-protocol}

The metrics are those of Section~\ref{sec:method-withdrawals} and the equations cited there: the gap index and the risk-weight difference with percentile intervals over $\bootB$ resamples of the cohort years (seed 20261104, and seed \fiPopSeed{} for Fitch's $\fiPopYears$ populated years); the relative bias and the risk-weight error with their Monte Carlo standard errors over $R$ repetitions; the rejection rate with the Wilson interval; the exact coverage; and, for the theory, the standardised difference $z=(\bar{\hat p}-p^{\ast})/\operatorname{se}$ between a simulated mean and its closed form $p^{\ast}$. The committed computations of the case, with 200 repetitions per level of $k\in\{0,1,3,9\}$ and the seeds derived from the case seed 42, are first reproduced; the largest difference between the committed means and the reproduced means is $\reproMax$. The extension then runs $\repsNew$ repetitions per level for $k\in\{0,1,2,3,5,9,19\}$ on seeds of its own, with the four pooled estimators and their first-cohort versions, the duration estimator on the observed paths of the same seeds, and the transition tests over $\repsTests$ repetitions at $k\in\{0,9,19\}$. The theory is checked on a two-state chain with the PDs of grades B and CCC to C, 40 paired draws of $\thObligB$ and $\thObligC$ obligors respectively, five levels of $k$, and both scopes; the first-cohort and pooled scopes give $\thNcmpFirst$ and $\thNcmpPooled$ distinct comparisons ($\thNcmp$ entries, of which the first-cohort followed estimator and the reference read the same draws and have the same closed form, so that each such pair counts once). The coverage is enumerated on a grid of cohort sizes from 50 to $\gridMax$ obligors and at the cohort sizes of the generator. The extension takes about \minWd{} minutes of computation, the theory check about \minTheory{} minute and the duration estimator with the further tests about \minExtra{} minutes.

\subsection{Ablations planned}
\label{sec:design-ablation}

The protocol has four components and several design choices inside them. The ablation of Section~\ref{sec:ablation} removes one at a time and measures, on the same data, the quantity that the component made visible: the side-by-side definitions and their capital, the ladder of informativeness, the known-truth reference, the pooling of five cohorts against the first cohort alone, the pairing of the observed and latent draws, the exact enumeration of the coverage against its simulation, and the matrix tests as a screen. Each row of Table~\ref{tab:ablation} states the component, the variant without it, and the two values of the quantity. The first two rows read the committed artifacts of the real cohorts; the next four read the paired draws of the extension and use the seeds of the main experiment, so that the variant without the component is computed on the very datasets the component was computed on; the last two read the exact enumeration and the matrix tests. The ablation does not search over combinations, because the components do not share an input except the generator, and so no interaction between them is claimed.
